\documentclass[11pt,a4paper]{article}
\usepackage[utf8]{inputenc}
\usepackage[T1]{fontenc}
\usepackage[spanish]{babel}
\usepackage{hyperref}
\usepackage{geometry}
\geometry{margin=2.5cm}
\usepackage{parskip}

\title{089 toda vision tiene jugo y entregables}
\author{Alejandro Valenzuela Rivera\\\small Universidad Autónoma de Nuevo León}
\date{\today}
\begin{document}
\maketitle

\begin{quote}
\textbf{Sesgo declarado.} Este trabajo se sostiene desde el sesgo declarado de su autor, y no finge neutralidad. Lo que afirma se puede tumbar, y ese es el punto: el falsificador está escrito, no escondido.
\end{quote}

\section*{Resumen}
El jugo dice de donde viene; el entregable dice que se sostiene. Sin jugo, el entregable no tiene genealogia; sin entregable, el jugo no llega a nadie. Y el jugo no tiene estandar: es mutable, crece y se corrige.

\section{Por qué se sostiene}
Son dos folders obligatorios porque son las dos mitades de una vision. Y el reparto es lo que permite que el jugo crezca sin frenar la entrega, y que la entrega no arrastre todo el jugo.

\section{De dónde viene}
La autoridad, y la vision del autor.

\section{El caso}
Ninguna: es la forma de la vision.

\section{Método}
Una vision que deja algo que se puede usar, ademas de lo que ensena.

\section{Límites, y qué lo tumba}
Una vision que solo explica, y no deja nada que se pueda correr.

\section{Lo que queda abierto}
\textquestiondown{}Que le falta a una vision que solo tiene jugo, o solo tiene entregable?

\section*{Cómo se leyó esta tesis}
Esto se escribió en el repositorio man-made-empathy, del que este candado es una pieza: candados/06 del caso/089 toda vision tiene jugo y entregables.md. La máquina que lo compuso lo hizo desde este mismo archivo, y no agregó nada que no estuviera escrito. Si algo de esto es falso, se dice en el mismo sitio, y se corrige con fecha.

\end{document}
